% ===================================================================
%  LE VOLUME COMMENCE ET FINIT A LA COUVERTURE IMPRIMEE, papier vert
%  (feuillets 3 et 4 en tete, 237 et 238 en queue). La RELIURE n'est
%  pas transcrite : les feuillets 1 et 240 sont les deux plats d'un
%  cartonnage d'editeur en percaline brune a titre dore, les feuillets
%  2 et 239 les gardes qui les doublent. Ils ont ete releves — le
%  premier plat porte trois lignes dorees et un filet maigre, le second
%  est nu — mais ils appartiennent a l'exemplaire, non a l'edition, et
%  la transcription s'arrete a ce que l'imprimeur a compose.
%  Le papier de la couverture se reconnait a la mesure : RVB median
%  177/163/139 aux feuillets 3 et 4, 177/169/152 et 176/175/161 aux
%  237 et 238, contre 237/210/184 au papier du corps d'ouvrage.
% ===================================================================

% ===================================================================
%  FEUILLET 3 — LA COUVERTURE IMPRIMEE (premiere page de couverture).
%  Le volume porte donc a la fois une reliure d'editeur (feuillets 1 et
%  240) ET la couverture imprimee, conservee a la reliure.
%
%  CALIBRATION : le BORD SUPERIEUR DU PAPIER, pris au milieu de la
%  rampe d'ombre du scan — y = 38 px, +- 4 px — puis la constante
%  C = 5,84 mm du feuillet 7. Incertitude totale +- 1,6 mm sur la
%  position d'ensemble ; les ecarts internes, eux, sont mesures au
%  pixel. Desinclinaison -1,60 degre.
%
%  LE BORD GAUCHE EST ROGNE PAR LA NUMERISATION : a x=0 la luminance
%  vaut 213 et DECROIT vers l'interieur — il n'y a pas de rampe
%  d'ombre, donc le papier continue hors champ. Il manque environ 33 px
%  = 2,8 mm a gauche, et la largeur mesurable n'est que de 108 mm pour
%  les 116 du volume. L'AXE DE CENTRAGE EST DONC ESTIME, NON MESURE.
%
%  « LINGUO INTERNACIONA » DEBORDE LA JUSTIFICATION : 1148 px pour les
%  1083 du volume, soit 5,6 mm de trop. C'est le meme debordement que
%  la page de titre, et \VUcentreA le laisse deborder au lieu de le
%  replier, la ligne etant posee dans une \hbox et non dans un alinea.
%
%  LA VIGNETTE est l'embleme de l'Ido — une etoile a six branches
%  portant IDO, cerclee de la legende LINGUO INTERNACIONA DI LA
%  DELEGITARO. Comme le portrait du feuillet 7 et le fleuron du folio
%  232, elle est DECOUPEE DU SCAN : rot[1207:1456, 482:723], 241 x 249
%  px = 20,40 x 21,08 mm.
%
%  LES DEUX DERNIERES LIGNES N'EN FONT QU'UNE : « Imprimita en
%  Luxemburg. » et « Preco : 2 Suisa Franki. » partagent la ligne de
%  base y=1893, l'une au fer a gauche, l'autre au fer a droite. D'ou
%  \VUcaseA, qui pose les deux sur une mesure donnee.
%
%  DEUX GRAISSES NON TRANCHEES, et dites comme telles :
%   -- « DI LA » : rapport fut/capitale 0,250, aussi haut que le IDO
%      gras, mais sur une capitale de 20 px ou un pixel de fut deplace
%      le rapport de 0,05. Agrandi x6, c'est un romain ordinaire un peu
%      engraisse par l'etalement. Compose en romain, non prouve.
%   -- « 1925 » : 0,174, entre le romain (0,154) et le demi-gras
%      (0,192) de la meme page. Compose en romain, sans certitude.
% ===================================================================
\begin{VUpage}[3]{}
\VUcentreA{16.25mm}{10.63pt}{109}{\VUgras{L. de BEAUFRONT}}
\VUfiletA[1.92pt]{23.79mm}{12.70mm}
\VUcentreA{37.51mm}{19.13pt}{112}{\VUetroit{0.668}{KOMPLETA GRAMATIKO DETALOZA}}
\VUcentreA{50.29mm}{7.09pt}{276}{DI LA}
\VUcentreA{61.04mm}{23.74pt}{29}{LINGUO INTERNACIONA}
\VUcentreA{76.54mm}{41.09pt}{4}{\VUgras{IDO}}
\VUimageA{105.06mm}{20.40mm}{ornamenti/vinieto-3}
\VUcentreA{147.32mm}{8.86pt}{0}{\VUgras{Editerio : MEIER-HEUCKE}}
\VUcentreA{152.23mm}{9.21pt}{0}{\VUgras{Esch-Alzette} \VUgras{\textit{(Luxemburgia)}}}
\VUfiletA[0.96pt]{156.97mm}{2.62mm}
\VUcentreA{160.10mm}{8.15pt}{80}{1925}
\VUcaseA{164.84mm}{102.53mm}{9.21pt}{Imprimita en Luxemburg.}{Preco : 2 Suisa Franki.}
\end{VUpage}

% ===================================================================
%  FEUILLET 4 — LA DEUXIEME PAGE DE COUVERTURE : « QUO ESAS IDO? »
%  Ce n'est pas un verso blanc, comme le disait le LISEZ-MOI : c'est
%  une page PLEINE, 45 lignes, justification 1080 px = 91,44 mm, soit
%  a 0,25 mm pres celle du volume.
%
%  LE PAS DES LIGNES EST DE 33 px (2,79 mm), NON 41 : c'est
%  l'interlignage des NOTES du volume, et exactement celui des
%  feuillets 237-238, les annonces de l'editeur, deja composes avec
%  \VUpageSerre{8.26pt}{7.95pt}. Hauteur de capitale du corps 23-24 px.
%  Le renfoncement d'alinea vaut 34 px = 2,88 mm, non les 4,6 mm de
%  \VUalinea : il est pose dans la page.
%
%  DEUX BLANCS D'ALINEA DISTINCTS, et c'est la mesure la plus utile de
%  la page : dans la section haute, exces 7-8 px = 1,69 a 1,93 pt,
%  c'est-a-dire \VUblancAlinea (1,86 pt) ; dans la section « Opinioni »,
%  exces 12 px = 2,89 pt.
%
%  CONVENTION D'ENRICHISSEMENT INVERSEE, et conservee telle quelle :
%  la section 3.3 du LISEZ-MOI donne gras = forme ido en vedette,
%  italique = mot d'une autre langue. Ici les mots IDO en vedette —
%  Komitato, Delegitaro, Progreso, un, Ido — sont composes en
%  ITALIQUE. C'est coherent : cette page n'est pas du corps du volume,
%  c'est un texte d'editeur.
%
%  « 1,250 » avec la virgule des milliers est bien ce qui est imprime.
%
%  LA PAGE S'ARRETE AU MILIEU D'UNE PHRASE — « ...ta principi aplikesis
%  kun rigoro » — et la suite n'est nulle part dans le scan : la
%  troisieme page de couverture n'a pas ete reliee. La derniere ligne
%  porte donc \nl et non \parplein : ce n'est pas une fin d'alinea.
% ===================================================================
\begin{VUpage}[4]{}
\VUcentreA{16.04mm}{14.17pt}{45}{\VUgras{QUO ESAS IDO?}}
\VUfiletA[1.44pt]{23.66mm}{8.13mm}
\VUcentreA{118.57mm}{9.92pt}{73}{\VUgras{Opinioni pri Ido.}}
\vspace*{4.61mm}
\VUpageSerre{8.26pt}{7.95pt}
\setlength{\parindent}{2.88mm}

Ido esas l'idiomo internaciona, helpanta e neutra, qua rezultas\nl
de la selekto e labori dil « \textit{Délégation} » (delegitaro).

\VUblanc{1.86pt}\textit{La Délégation pour l'adoption d'une langue auxiliaire interna}\cc
\textit{tionale}, fondita en 1901 (januaro), esis ensemblo de personi delegita\nl
da 310 kongresi o societi omnalanda ed omnaskopa, por examenar la\nl
multa projeti o sistemi di linguo internaciona, e por adoptar la\nl
maxim bona. Til 1907 ol recevabis 1,250 aprobanta signaturi de\nl
membri di Akademii od Universitati, qui tale montris sua intereso\nl
al entreprezo.

\VUblanc{1.86pt}Lore ol elektis (junio 1907) internaciona \textit{Komitato} di ciencisti e\nl
linguisti partikulare kompetenta, qui, pos exameno di omna projeti\nl
ancien o recenta di Linguo internaciona ante e dum 18 kunsidi,\nl
adoptis fine Esperanto kun diversa modifiki. Ta modifiki studiita,\nl
kompletigita per laboro komuna, publik ed internaciona, de 1908 til\nl
1914, en la revuo \textit{Progreso}, havis kom rezultanto la nuna « Ido ».\nl
Ica konseque ne esas, quale Esperanto, la simpla produkturo di \textit{un}\nl
homo.

\VUblanc{1.86pt}Esas notenda, ke la Komitato dil \textit{Delegitaro} kontenis Esperantisti\nl
e mem la Prezidero ipsa di la « Lingva Komitato » esperantista.\nl
Or la decido adoptanta Esperanto « \textit{kun la rezervo di ula modifiki} »,\nl
esis unanima. Se, pose, skismo eventis, la responso ne esas imputebla\nl
a l'Idistaro, ma ad ilti qui obliviis promisi, voti o signaturo.

\VUblanc{1.86pt}La Komitato dil Delegitaro (la maxim kompetenta e maxim sen\cc
partisa qua ultempe kunvenis pri linguo internaciona) tale judi\cc
kesas da S\textsuperscript{ro} \textsc{Gaston Moch}, qua supleis en lu rektoro \textsc{Emile Boirac},\nl
tre eminenta chefo Esperantista : « La lingui Germana e Franca\nl
havis ibe tri reprezentanti, la Angla du, e la lingui Dana, Hispana,\nl
Greka e Hungariana, un; la maxim diversa konocaji reprezentesis,\nl
ed on remarkis en lu nome quar reputata filologi, S\textsuperscript{ri} \textsc{Baudouin}\nl
\textsc{de Courtenay}, \textsc{Jespersen}, \textsc{Lambros} e \textsc{Schuchardt}. »

\VUblanc{18.32pt}« \textit{La gramatiko di Ido satisfacas plu bone la postuli di linguo}\nl
\textit{internaciona kam la gramatiko di Esperanto...} » (A. \textsc{Meillet}, pro\cc
fesoro pri linguistiko en Collège de France, \textit{Revue Critique}, 11 di\nl
marto 1911.)

\VUblanc{2.89pt}« Esas eroro nepardonebla institucar, quale agis Esperanto,\nl
distingo pri l'akuzativo e la nominativo, distingo qua jenos omna\nl
individui di linguo romanal od Angla, e qua esas neutila a la ceteri. »\nl
(La sam autoro, \textit{Revue Critique}, 11 di marto 1911.)

\VUblanc{2.89pt}« Cetere, esas facila procedar plu logikoze e, konseque plu satis\cc
facante e klare, en la vortifado kam agas lo Esperanto.

\VUblanc{2.89pt}« Icon montris la kreinti di \textit{Ido}, linguo fondita sur la sama\nl
principi kam Esperanto, ma ube ta principi aplikesis kun rigoro\nl
\end{VUpage}

% ===================================================================
%  Feuillets 5 et 6 — pieces liminaires (folios 1 et 2)
%
%  Pages d'apparat : chaque ligne est placee a son ORDONNEE MESUREE et
%  occupe une hauteur nulle ; aucune ligne n'en deplace une autre.
%  Mesures par outils/apparat.py (feuillet 5) :
%    ligne       y (px)  y (mm)  cap (mm)  largeur (mm)  corps    interl.
%    KOMPLETA...    754   63,84    2,540      60,20      10,63pt      2
%    DI LA          919   77,79    1,439       7,03       6,02pt    224
%    LINGUO...     1063   89,98    3,852      89,15      16,12pt    117
% ===================================================================

% --- Feuillet 5 (folio 1) : faux-titre -----------------------------
\begin{VUpage}[5]{}
\VUcentreA{63.84mm}{10.63pt}{-5}{KOMPLETA GRAMATIKO DETALOZA}
\VUcentreA{77.79mm}{6.02pt}{264}{DI LA}
\VUcentreA{89.98mm}{16.12pt}{102}{LINGUO INTERNACIONA IDO}
\end{VUpage}

% --- Feuillet 6 (folio 2) : blanc ----------------------------------
\begin{VUpage}[6]{}
\end{VUpage}

% -------------------------------------------------------------------
%  FEUILLET 7 (folio 3) — LA PAGE DE TITRE, et la planche.
%  Dix elements, tous d'apparat : pas une ligne de texte courant.
%  Le releveur les a trouves centres sur un meme axe optique a moins de
%  deux pixels pres, et a montre que le detecteur en avait perdu deux
%  (LINGUO INTERNACIONA IDO, dont les capitales de 77 px depassent le
%  filtre de hauteur, et DA, large de 41 px) tout en en inventant six,
%  qui sont des artefacts de la trame de la similigravure.
%
%  LA PAGE N'A NI FOLIO NI LIGNE DE CORPS : la regle ordinaire — poser
%  l'ordonnee relativement a la premiere ligne de corps — n'a rien a
%  quoi s'accrocher. Le premier etat de cette page etait donc pose en
%  FLUX, avec un blanc de tete choisi a l'oeil. Il etait faux : la
%  comparaison du rendu au fac-simile a montre que CHAQUE blanc y etait
%  trop court, la page se comprimant de 139 px du haut vers le bas.
%
%  Elle est reprise ici en ordonnees ABSOLUES, mesurees sur le scan
%  desincline de 0,20 degre :
%    ordonnee = (haut_des_capitales - bord_du_papier) x 25,4/300 + C
%  ou C = 5,84 mm est la part de la marge de tete que la numerisation
%  rogne. C n'est pas un chiffre libre : il est mesure sur 33 pages de
%  texte ou la premiere ligne de corps tombe, par construction, a
%  \VUmargeSup ; C = 24,30 - (y_ligne - y_papier) x 25,4/300. La
%  mediane vaut 5,84 mm, l'ecart-type 1,5 mm hors trois pages de titre.
%  C'EST LA SEULE GRANDEUR APPROCHEE DE CETTE PAGE, et elle porte sur
%  la position d'ensemble du bloc, non sur ses ecarts internes, qui
%  sont mesures exactement.
%
%  Le haut des capitales du fac-simile est corrige de l'etalement de
%  l'encre : (hauteur_fac - hauteur_compose)/2, soit 0,5 a 6 px.
%
%  LA LIGNE « LINGUO INTERNACIONA IDO » DEBORDE LA JUSTIFICATION :
%  1164 px pour les 1083 du volume. C'est mesure trois fois, ce n'est
%  pas la tranche du feuillet voisin. \VUcentreA la pose dans une
%  \hbox et non dans un paragraphe : elle deborde comme au fac-simile
%  au lieu de se replier.
% -------------------------------------------------------------------
\begin{VUpage}[7]{}
\VUcentreA{10.71mm}{16.60pt}{-68}{KOMPLETA GRAMATIKO DETALOZA}
\VUcentreA{20.36mm}{6.30pt}{362}{DI LA}
\VUcentreA{27.94mm}{24.10pt}{-79}{LINGUO INTERNACIONA IDO}
\VUcentreA{40.00mm}{5.90pt}{242}{DA}
\VUcentreA{45.17mm}{8.80pt}{282}{\VUgras{L. de BEAUFRONT}}
\VUcentreA{49.06mm}{5.90pt}{222}{PREZIDERO DI LA FRANC-IDISTA SOCIETO}
\VUplancheA{55.71mm}
\VUcentreA{148.76mm}{5.90pt}{124}{Markezo L. \textsc{de Beaufront}, precipua autoro di \textit{Ido}}
\VUcentreA{157.31mm}{7.60pt}{148}{Imprimerie SOLIMPA, Luxembourg et Esch/Alzette}
\VUcentreA{160.99mm}{7.60pt}{-21}{1925}
\end{VUpage}

% ===================================================================
%  FEUILLET 8 (folio 4) — LE VERSO DE LA PAGE DE TITRE : la KONSTATO
%  de l'Akademio Idista, et sa signature.
%
%  Les 266 473 px d'encre annonces au seuil 140 sont trompeurs : 54 934
%  seulement sont de l'encre typographique, le reste est le fond noir
%  du scanner autour du papier (papier : y 37-2109, x 34-1310).
%
%  LA PAGE EST DELIBEREMENT AEREE : pas des lignes de base 57,667 px =
%  13,89 pt, contre 9,99 pt au corps courant, a CORPS INCHANGE (hauteur
%  d'x 20-21 px, celle du feuillet 10). D'ou \VUinterlignePage.
%
%  CALIBRATION : la page n'a ni folio ni ligne de corps a \VUmargeSup —
%  son corps commence a 54,27 mm parce qu'un titre et une marge de tete
%  profonde le precedent. C'est donc la calibration du feuillet 7 qui
%  sert : (y - bord_du_papier) x 25,4/300 + C, bord a y=37, C=5,84 mm.
%  RESERVE DU RELEVEUR : son propre echantillon local (feuillets 10, 12
%  et 16, les seuls voisins ou le detecteur rende folio ET premiere
%  ligne) donne C = 6,27 mm. Le bloc entier peut donc etre 0,43 mm trop
%  haut. Trois echantillons contre trente-trois : la constante du
%  projet est gardee, et l'ecart est dit.
%
%  LE BLOC DE SIGNATURE N'EST PAS CENTRE SUR LA JUSTIFICATION : ses
%  quatre lignes ont pour milieux 900,5 / 901,5 / 900,5 / 901,0 px pour
%  un milieu de justification a 706,5 — un axe propre, deplace de
%  16,46 mm a droite, constant a un pixel pres. D'ou \VUcentreDA.
%  La date, elle, est calee a GAUCHE, a 80 px = 6,77 mm de la marge.
%
%  GRAISSE : le titre est demi-gras (fut du K 5,94 px contre 4,78 au K
%  romain de « Komitato »), mais sa hauteur de capitale, 29 px, est
%  celle du CORPS COURANT et non celle d'un titre de section (44 px au
%  feuillet 9). Dans le corps, seul le titre cite « Kompleta Gramatiko
%  detaloza di la linguo internaciona Ido » est gras : futs 4,39 a 4,62
%  px contre 3,64 a 3,95 au romain de la meme ligne. Le mot « naciona »
%  etait le seul cas douteux — densite de gras, fut de romain : agrandi
%  x6 il est de la meme graisse que « quala » qui le suit, sa densite
%  etant gonflee par l'absence d'ascendante et de descendante dans la
%  bande de mesure. Il est donc ROMAIN.
%
%  PARTICULARITE CONSERVEE : le meme syntagme parait deux fois de
%  suite avec deux traitements opposes — en gras et bas de casse comme
%  TITRE D'OUVRAGE, puis en romain a capitale initiale comme NOM DE
%  LANGUE. Ce n'est pas une coquille, c'est la distinction voulue.
%  Le renfoncement d'alinea de cette page vaut 42 px = 3,56 mm, et non
%  les 4,6 mm de \VUalinea ; il n'est pas force.
% ===================================================================
\begin{VUpage}[8]{}
\VUinterlignePage{13.89pt}
\VUcentreA{44.53mm}{10.06pt}{119}{\VUgras{\VUetroit{0.793}{KONSTATO.}}}
\VUcentreDA{91.52mm}{16.46mm}{10.2pt}{-13}{\textit{La prezidanto :}}
\VUcentreDA{96.43mm}{16.46mm}{10.2pt}{-66}{F. \textsc{Schneeberger.}}
\VUcentreDA{102.70mm}{16.46mm}{10.2pt}{-13}{\textit{La vice-sekretario :}}
\VUcentreDA{107.61mm}{16.46mm}{10.2pt}{-86}{J. \textsc{Guignon.}}
{\renewcommand{\VUsignatureX}{6.77mm}%
 \VUsignature{117.43mm}{\fontsize{10.2pt}{10.2pt}\selectfont\textls[-16]{\textit{15 julio 1925.}}}}
\vspace*{29.97mm}

L'Akademio Idista konstatas, ke la linguala formo pri\cc
zentita en ica \VUgras{Kompleta Gramatiko detaloza di la linguo}\nl
\VUgras{internaciona Ido} esas la oficala formo di la Linguo inter\cc
naciona \VUgras{Ido}, quala ol rezultis de la labori di la Linguala\nl
Komitato di la Delegitaro por l'adopto di Linguo Interna\cc
ciona, di lua konstanta komisitaro e la decidi dil Akademio\nl
Idista.
\end{VUpage}

% -------------------------------------------------------------------
%  FEUILLETS 9 a 12 (folios 5 a 8) — L'AVERTO.
%  Quatre pages relevees et verifiees separement, une par releveur.
%  Aucun gras dans le corps ni dans les notes des quatre pages :
%  l'italique y est partout plus MAIGRE que le romain voisin, mesure
%  fut par fut contre un mot romain de la MEME ligne.
% -------------------------------------------------------------------

% --- Feuillet 9 (folio 5) : le titre AVERTO ------------------------
%  AUCUN FOLIO IMPRIME sur cette page : cherche au seuil 165 sur toute
%  la largeur, a \VUfolioY comme au pied. L'argument reste vide.
%  AUCUN FILET sous le titre : cherche entre y=471 et y=557 aux seuils
%  150, 170, 185 et 200 — rien que le grain du papier.
%  Le titre « Averto » est en BAS DE CASSE demi-gras, capitale initiale,
%  hauteur de capitale 44 px (3,73 mm) : c'est le seul titre du volume
%  qui ne soit pas en capitales. Fut 8 px pour une hauteur d'x de 27 px
%  (0,296) contre 4 px pour 21 px au romain de la page (0,190).
%  Son appel « (1) » est MAIGRE, au corps du texte, remonte de 13 px.
%  Filet de note : y=1613, 213 px, cale sur la marge (x0=90 pour une
%  premiere ligne de note a x0=89) — 123,36 mm.
%  ATTENTION : l'ordonnee est calculee sur le TITRE, non sur une ligne
%  de corps. Recalculee sur la premiere ligne de corps elle donnerait
%  113,54 mm, soit 10 mm trop haut.
%  Le bloc de la note (2) est pose 23 px (2 mm) plus a droite que celui
%  de la note (1) et que le filet — verifie sur l'image BRUTE, ce n'est
%  pas un artefact de desinclinaison. Signale, non reproduit.
% -------------------------------------------------------------------
\begin{VUpage}[9]{}
\VUnotes{122.78mm}{%
(1) L'unesma parto di ca averto esas nur la vortopa tradukuro\nl
di olta qua aparis France en 1908, lor l'edito di \textit{Grammaire complète}.

(2) \textit{La Délégation pour l'adoption d'une Langue auxiliaire inter}\cc
\textit{nationale}, fondita en 1901 (januaro), recevabis til 1907 la adhero\nl
di 310 societi omnalanda, e la aprobo di 1250 membri di Akademii\nl
ed Universitati. Ol elektis internaciona Komitato di ciencisti e lin\cc
guisti, qua pos exameno di omna projeti anciena e nova di mondo\cc
linguo, facis la decido quan on raportas hike.%
}

\VUtitre{15.66pt}{-8}{\VUetroit{0.822}{Averto}\VUtitreAppel[1.83mm]{(1)}}
\VUsaut{5.0mm}

Ica verko esas ri-imprimuro grande modifikita di la\nl
\textit{Grammaire complète} qua, imprimita nur ye 200 exempleri\nl
e ne publikigita, livresis al dispono dil Komitato di la Dele\cc
gitaro en la 15-ma dio di oktobro 1907, kun \textit{Exercaro} e\nl
\textit{Vortaro} specimena, sub la pseudonimo « Ido » (2).

On savas, ke, pos 18 kunsidi, ek qui 5 konsakresis al dis\cc
kuto komparanta di ta laboruro e di Esperanto, unanime\nl
la Komitato decidis « adoptar Esperanto principe... kun la\nl
rezervo di ula modifiki exekutenda dal permananta Komi\cc
sitaro, segun la sinso determinita dal konkluzi di la Raporto\nl
dil sekretarii e dal projeto da Ido ». Ja la Komisitaro indi\cc
kabis ula modifiki facenda en ta projeto. Sualatere la Komi\cc
sitaro permananta submisis lu a detaloza revizo, egardante\nl
plura emendi propozita; e la rezultajo di ta laboro men\cc
cionesis en Raporto publikigita en l'unesma numero di la\nl
revuo \textit{Progreso} (marto 1908).

La \textit{Grammaire complète} quan ni publikigas prezente esas\nl
strikte konforma al decidi di la Komitato e dil Komisitaro\nl
permananta; do ol ne plus esas, same kam la linguo ipsa,\nl
verko pure individuala; ol esas la final ed oficala rezultajo\nl
di la deliberi dil Komitato elektita reguloze dal \textit{Delegitaro},\nl
pos sep yari de propagado por l'ideo di linguo internaciona\nl
en la maxim diversa landi e medii. Ni kun fido prizentas\nl
lu al publiko pro la alta kompetenteso di la ciencozi qui\nl
partoprenis en lua kompozo.
\end{VUpage}

% --- Feuillet 10 (folio 6) -----------------------------------------
%  Filet de note : y=1527, 213 px, x0=196 pour une marge a 195-199.
%  Le detecteur a bien pris la premiere ligne de CORPS (y=247) et non
%  le folio (y=171, qui aurait donne 139,10 mm).
%  Les dix lignes de note sortaient du profil SOUDEES en quatre plages ;
%  separees aux lignes de base, pas d'apres le profil.
% -------------------------------------------------------------------
\begin{VUpage}[10]{6}
\VUnotes{133.14mm}{%
(1) Kontante segun la maniero dil Esperantisti, ni povus dicar,\nl
ke Ido havas nur 10 e mem 9 reguli, ne min kompleta. (Videz\nl
« Lexique-manuel », p. V-VIII.) --- \textit{Noto adjuntita.}

(2) En la realeso la gramatiko \textit{elementala} di Esperanto kontenas\nl
adminime 64 reguli e plu kam 10 ecepti. (Videz la konto detalizita\nl
ye la fino dil broshuro « \textit{La science et l'Esperanto} ».)

(3) Hike ni abandonas la Franca texto, qua deskriptas l'aranjeso\nl
di \textit{Grammaire complète}, e ni substitucas ad olu to, quo propre\nl
koncernas la « \textit{Kompleta Gramatiko detaloza} » quan ni prizentas\nl
ad omna amiki di helpolinguo.%
}

La kritiki, di qui ca gramatiko esos kredeble l'objekto,\nl
ne falios opozar ad olu la mikrega gramatiko primitiva di\nl
Esperanto, kun lua dek e sis reguli tante laudata (1). Ma,\nl
sen diskutar pri la numerizo di ta reguli, o plu \textit{vere para}\cc
\textit{grafi}, di qui un sola kontenas la konjugo (minus la formi\nl
kompozita per qui on mustis kompletigar lu pose), on darfas\nl
asertar sen ul paradoxo ke, se Esperanto havabus origine\nl
plu multa reguli, ol nun posedus quanto min granda de\nl
oli (2). Nam pro ta gramatiko vere tro kurta e nesuficanta,\nl
fakte a l'uzado Dro Zamenhof komisis la sorgo fixigar la\nl
formi di la linguo, ed ica uzado genitis multega partikularaji\nl
ed anomalaji, quin la gramatikisti di Esperanto kolektas e\nl
pie mencionas en lia verki, quale se traktesus linguo natu\cc
rala. E singla de ta partikularaji, singla de ta anomalaji\nl
efektigas un « regulo » quan la novici esas obligata lernar.\nl
Do esas tre vera dicar, ke se Esperanto havabus origine sat\nl
multa reguli generala e preciza, ol ne havus nun tro multa\nl
partikulara reguli, fondita sur « uzi » disparata e kontre\cc
logika. Pro to on obligesis kompozar fine grosa gramatiki,\nl
de qui un prizentas plu kam 400 paragrafi. A ca verki ya\nl
on devas equitatoze komparar nia \textit{Grammaire complète}. Lore\nl
on ne povos neagnoskar la supera simpleso e regulozeso di\nl
la linguo quan ni prizentas (3).

Certe ni povabus kondensar e multe plukurtigar ica verko,\nl
nome ne enduktar en lu kozi advere utila ma ne necesa, e\nl
qui ne esas strikte gramatikala. Ma lore ni agabus kontre\nl
nia intenci e l'explicita deziro di multa samideani. Cetere\nl
on dicernos facile la reguli lernenda del kozi nur lektebla\nl
o lektinda.

Kompreneble, en la gramatiko sancionita dal permananta
\parplein
\end{VUpage}

% --- Feuillet 11 (folio 7) -----------------------------------------
%  Filet de note : y=1726, 212 px, x0=124 pour un corps a x0=121-126.
%  Aucun filet centre, aucun ornement : rien entre 224 et 277, entre
%  1686 et 1726, ni entre 1916 et 2122.
% -------------------------------------------------------------------
\begin{VUpage}[11]{7}
\VUnotes{147.07mm}{%
(1) Ica vorti di nia unesma averto divenas partikulare vera pri\nl
la \textit{Kompleta Gramatiko detaloza,} en qua la multa pruntaji ek \textit{Pro}\cc
\textit{greso} e citaji di Sro \textsc{Couturat} posibligos ta profunda studiado.\nl
Danke oli, ni darfas dicar : mortinta, il parolas ankore, la granda\nl
kalumniato.%
}

\VUcontinue Komisitaro ni permisis a ni chanjar od adjuntar, pri la\nl
reguli, nur la kelka punti chanjita od adjuntita dal Ido-\nl
Akademio, legitima e yurizita sucedinto di olta. Ma la formo\nl
ed aranjeso necese diferas de olti, quin havis la verko pri\cc
mitiva, haste facita en tri monati e nur destinita al membri\nl
dil Komitato dil \textit{Delegitaro,} por utiligo eventuala. Certe la\nl
permananta Komisitaro ja emendabis ed igabis lu plu bona,\nl
ma vizante la linguo ipsa e ne olua docado specale. Or la\nl
publiko tre diferas de komitato di ciencisti, por qua multa\nl
kozi bezonis nek detali nek justifiko. Cetere la tempo mankis\nl
al Komisitaro por rilaborar detale la verko, nam adversi\nl
ed amiki ja esis incitanta a la propagado, iti kun intenco\nl
maligna, ici en sua fervoro difuzor lo plu bona. Rezulte, la\nl
primitiva gramatiko bezonis ankore plu kompleta riaranjo\nl
ed expliki plu longa. Ni esforcis en ica satisfacar ta duopla\nl
bezono e preferis donar tro multe kam ne sufice.

On permisez ad ni atraktar la atenco a kozo quan multi\nl
ne vidas sat bone. Pri la helpolinguo existas facileso trom\cc
panta, qua vizas nur o precipue \textit{lerno rapida} per nesuficanta\nl
reguli e nesuficanta radiki o vorti. Nu, ica facileso \textit{di lerno}\nl
ne povas konkordar e koincidar kun la facileso \textit{di apliko,} nek\nl
kun l'exakta e klara expresado dil pensi. Or on lernas la\nl
L. I. dum kelka hori, kelka jorni, ed on aplikas lu dum yari.\nl
Altraparte, quon valoras facileso e simpleso qui reale ne\nl
atingas la skopo? Mezurar la supereso, pri la helpolinguo,\nl
segun la minmulteso di la reguli e dil vorti, forsan esas\nl
habila trompilo por la maso, ma eroro deplorinda por omni,\nl
eroro qua produktas en la aplikado detrimenti maxim grava.\nl
Ido evitis fortunoze ta rifo mortigiva, kontre qua mult altri\nl
ja venis ed ankore venos ruptar su.

Nun ni finos, durante la traduko quan ni interruptis :

« Cetere, ica \textit{Kompleta Gramatiko} destinesas precipue al\nl
personi qui volas profundigar la studiado di la linguo (1).\nl
Ol esforcas solvar omna desfacilaji quin on povas renkontrar\parplein
\end{VUpage}

% --- Feuillet 12 (folio 8) : la fin de l'AVERTO --------------------
%  AUCUN BLOC DE NOTES : « filet non trouve » est ici la bonne reponse.
%  Le seul trait de la page part de x=648 et mesure 197 px : c'est
%  l'ORNEMENT DE FIN, centre (milieu 746 pour une justification centree
%  a 755), pose a 149,01 mm.
%  DEUX SIGNATURES au fer a droite, mais reculees de 41 et 43 px
%  (3,47 et 3,64 mm) : le meme cadratin aux deux — d'ou \VUauFer.
%  « L. B. » AVAIT ECHAPPE AU DETECTEUR : 769 px d'encre contre 6882
%  pour une ligne pleine, soit 0,112 pour un seuil pose a 0,12. La page
%  porte 31 rangees d'encre, non 30.
%  « Sro » est en PETITES CAPITALES sur cette page, non en superieure
%  comme partout ailleurs dans le volume (temoin : feuillet 123).
%  « Chandelliez » est imprime avec deux L : conserve.
% -------------------------------------------------------------------
\begin{VUpage}[12]{8}
\VUornement[16.68mm]{149.01mm}

\VUcontinue en l'uzado di ca linguo, e qui venas ne de ol ipsa, ma del\nl
neregulozaji e del idiotismi di nia lingui. Ta desfacilaji ren\cc
kontresas en omna traduko, e mem duopligita, kande on\nl
traktas naturala lingui, nam generale on esas obligata\nl
expresar l'idiotismi di un linguo per l'idiotismi di altra. Ma\nl
precize ta desfacilaji facas ek la L. I. bonega intelektal\nl
exerco, analoga (egardite omna proporciono) a la studio di\nl
la Latina o di la stranjera lingui :

Nam li kustumigas la spirito analizar la penso e liberigar\nl
ol del vesto plu o min trompiva, per qua nia lingui vesti\cc
zachas olu. Ma, vice quik impozar a lu nova travestio, la L. I.\nl
tendencas furnisar a lu l'expresuro maxim direta, maxim\nl
logikala e maxim klara. E, malgre l'opozata dici, logiko\nl
ya vere igas la L. I. supera a nia lingui « naturala ». Logiko\nl
ya sekurigas ad olu la palmo en lua aplikeso a la cienci;\nl
logiko donas ad olu l'astonanta facileso, kun qua omna\nl
populi ed omna klasi sociala povas komprenar olu, aquirar\nl
olu ed uzar olu.

\VUblanc{1.95mm}
\VUauFer{3.47mm}{\textsc{L. de Beaufront.}}
\VUblanc{3.13mm}

Me falius devo di yusteso e gratitudo, se me ne dankus\nl
hike \textsc{Sro Camille Chandelliez}, ex-docero dil grupo Espe\cc
rantista en \textit{Troyes}, Francia. Konvertita ad Ido per plur\cc
monata lekto e studio di \textit{Progreso}, il ofris a me sua serchi e\nl
noti, sua remarki e kritiki sagaca. Il helpis me tre efike\nl
klarigar ula punti, nome l'adjektivi-pronomi nedefinita. Ad\nl
ilu komplete debesas l'apendico pri la puntizado, la letri e\nl
adresi, segun \textit{Progreso}. Il esis mea unesma lektero e kon\cc
stanta kritikero. Ica gramatiko debas multo a lua helpado.\nl
Ton me joyas dicar publike e gratitudoze.

\VUblanc{1.95mm}
\VUauFer{3.64mm}{\textsc{L. B.}}
\end{VUpage}

% ===================================================================
%  FEUILLET 13 (folio 9) — FAUX-TITRE DE LA PREMIERE PARTIE.
%  Le precedent exact est le feuillet 119, faux-titre de la deuxieme
%  partie : memes corps (14,17 pt a la grande ligne dans les deux cas),
%  memes intervalles a 0,25 mm pres, moins la ligne de sous-titre.
%
%  La page n'a ni folio ni ligne de corps. Calibration employee : celle
%  du feuillet 7 — (haut_des_capitales - bord_du_papier) x 25,4/300 + C,
%  avec C = 5,84 mm. Le releveur a RECALCULE C sur 34 pages a bord
%  visible : mediane 5,82 mm, ecart-type 1,15 — la constante du projet
%  est confirmee.
%  CONTRE-EPREUVE INDEPENDANTE : le feuillet 12 sort du MEME scan brut
%  (moitie gauche de la meme prise) et porte, lui, folio et texte. En
%  passant par sa premiere ligne de corps (y=275), la regle ordinaire
%  du volume donne 72,31 / 81,79 / 90,76 mm. Les deux calibrations
%  s'accordent a 0,21 mm. Ce sont les premieres qui sont retenues.
%
%  LES DEUX LIGNES SONT ROMAINES, non grasses : fut median 5 px pour
%  une capitale de 29 (0,172) et 7 px pour 40 (0,175), quand le
%  demi-gras compose en donne 10. Il n'y a aucun romain sur ces lignes
%  pour servir de temoin — elles sont seules sur la page — de sorte que
%  la comparaison porte sur le romain COMPOSE au meme corps. La methode
%  a ete controlee en la reappliquant au feuillet 119, ou elle redonne
%  les chiffres deja consignes.
%
%  Les interlettrages sont CORRIGES SUR LE COMPOSE, non pris de la
%  regle nue d'apparat.py, qui rendait 111 pour la grande ligne (2,5 mm
%  trop court). Les valeurs posees rendent 216 et 813 px, soit les
%  mesures du fac-simile au pixel pres.
%
%  RESERVE PORTANT SUR UNE AUTRE PAGE : le releveur soutient, mesures a
%  l'appui, que le FEUILLET 119 est pose 7,7 mm trop haut, ses
%  ordonnees etant les y x 25,4/300 bruts d'apparat.py, ce qui suppose
%  que la premiere rangee du recadrage soit le bord du papier — faux
%  sur cette prise-la. Voir le LISEZ-MOI, section des controles
%  ouverts : la reserve est enregistree, elle n'est pas encore
%  verifiee, et le feuillet 119 n'a pas ete touche.
% ===================================================================
\begin{VUpage}[13]{}
\VUcentreA{72.09mm}{10.27pt}{182}{I\textsuperscript{a} PARTO}
\VUfiletA{81.57mm}{8.30mm}
\VUcentreA{90.55mm}{14.17pt}{134}{MORFOLOGIO E SINTAXO}
\end{VUpage}

% --- Feuillet 14 (folio 10) : blanc --------------------------------
%  225 px d'encre au seuil 140 : la page est vierge.
\begin{VUpage}[14]{}
\end{VUpage}
